%----------------------------------------------------------------------------------------
%	PACKAGES AND OTHER DOCUMENT CONFIGURATIONS
%----------------------------------------------------------------------------------------
\DocumentMetadata{
pdfversion=2.0,  
lang=en-US,   
pdfstandard={ua-2, a-4f},
tagging = on, 
tagging-setup={math/setup={mathml-SE}} ,
}
\documentclass[12pt]{article}
\usepackage{amsmath}
\usepackage[extreme]{savetrees}
\usepackage[utf8]{inputenc}
\usepackage{newpxtext}
\usepackage{hyperref}
\usepackage{multicol}
\usepackage{physics}
\usepackage{siunitx}

%%%%%%%%% === Document Configuration === %%%%%%%%%%%%%%

\hypersetup{
pdftitle={MATH 340 -- Homework 9 Spring 2026},
pdfauthor={Ignacio Rojas},
pdfkeywords={ninth homework},%
}

\newcommand{\twobyone}[2]{\begin{pmatrix} %% 2 x 1 matrix
   #1 \\ #2 \end{pmatrix}}
\newcommand{\twobytwo}[4]{\begin{pmatrix} %% 2 x 2 matrix
   #1 & #2 \\ #3 & #4 \end{pmatrix}}
\newcommand{\threebyone}[3]{\begin{pmatrix} %% 3 x 1 matrix
   #1 \\ #2 \\ #3 \end{pmatrix}}
\newcommand{\threebythree}[9]{\begin{pmatrix} %% 3 x 3 matrix
   #1 & #2 & #3 \\ #4 & #5 & #6 \\ #7 & #8 & #9 \end{pmatrix}}

%----------------------------------------------------------------------------------------
%	ARTICLE CONTENTS
%----------------------------------------------------------------------------------------
\begin{document}

\begin{center}
   {\Large MATH 340 -- Homework 9,\quad Spring 2026}
\end{center}

Recall that you must hand in a subset of the problems for which deleting any problem makes the total score less than 10. The maximum possible score on this homework is 10 points. See the syllabus for scoring details. NOTE: All problems from Boyce, DiPrima, and Meade are from the \(11^{\text{th}}\) edition, which is available online: \href{https://www.math.colostate.edu/~liu/MATH340_S26Crd/Textbook_BoyceDiPrimaMeade_11ed.pdf}{https://www.math.colostate.edu/\(\sim\)liu/MATH340\_S26Crd/Textbook\_BoyceDiPrimaMeade\_11ed.pdf}.

\begin{enumerate}

   \item (1+) [2 points] Boyce, DiPrima, and Meade chapter 7 section 5 problem 9.

   \item (2-) [3 points] Boyce, DiPrima, and Meade chapter 7 section 5 problem 12.

   \item (2-) [3 points] Boyce, DiPrima, and Meade chapter 7 section 6 problem 19.
   
   \item Boyce, DiPrima, and Meade chapter 7 section 6 problem 25:
         \begin{enumerate}
            \item (1+) [2 points] parts a, b, c.
            \item (1+) [2 points] parts d, e, f.
         \end{enumerate}

   \item (1+) [2 points] Boyce, DiPrima, and Meade chapter 7 section 8 problem 4.

   \item (2-) [3 points] Boyce, DiPrima, and Meade chapter 7 section 8 problem 9.

   \item (1+) [2 points] Find the real-valued general solution of the system \(\vec{x}'=A\vec{x}\) where
         \[
            A=\threebythree{-5}{-6}{12}{3}{7}{-12}{0}{3}{-5}.
         \]

   \item (\textbf{E}) (1+) [2 points] Find the real-valued general solution of the system \(\vec{x}'=A\vec{x}\) where
         \[
            A=\threebythree{1}{-2}{0}{2}{1}{0}{0}{0}{-3}.
         \]

   \item (\textbf{E}) (1+) [2 points] Find the real-valued general solution of the system \(\vec{x}'=A\vec{x}\) where
         \[
            A=\threebythree{2}{-1}{1}{1}{2}{2}{0}{0}{3}.
         \]

   \item (2+) [4 points] Consider the system \(\vec{x}'=A\vec{x}\) where
         \[
            A=\frac12\threebythree{a+c}{-2b}{c-a}{b}{2a}{-b}{c-a}{2b}{a+c},
         \]
         and \(a,b,c\) are constants. Given that the eigenvalues of \(A\) are
         \[
            \lambda_1=a+ib,\quad\lambda_2=a-ib,\quad\text{and}\quad\lambda_3=c,
         \]
         Describe the behavior of a solution with initial condition \(\vec{x}(0)=(1,1,1)\) for different values of \(a,b,c\). Consider the possibilities where these constants are positive, negative, or zero.

   \item (1-) [1 point] The system of equations \(\vec{x}'=A\vec{x}\) with initial condition \(\vec{x}(0)=(1,0,1)\) and 
         \[
            A=\threebythree{1}{-3}{0}{0}{2}{0}{1}{1}{3}
         \]
         has its general solution given by
         \[
         \vec{x}(t)=c_1e^t\threebyone{2}{0}{1}+c_2e^{2t}\threebyone{3}{-1}{-2}+c_3e^{3t}\threebyone{0}{0}{1}.
         \]
      Find the particular solution with the given initial condition.

   \item (1+) [2 points] Consider a plane curve generated by the trace of a fixed point on a small circle of radius \(r= 1\) that rolls within a larger circle of radius \(R= 5\). Such a curve is called hypocycloid. The coordinates of this curve are expressed as
   \[
   \left\lbrace
   \begin{aligned}
      &x(t)=4\cos(t)+\cos(4t),\\
      &y(t)=4\sin(t)-\sin(4t).\\
   \end{aligned}
   \right.
   \]
   The coordinates of the hypocycloid satisfy the differential equations
   \[
   \left\lbrace
   \begin{aligned}
      &\ddot{x}-3\dot{y}+4x=0,\\
      &\ddot{y}+3\dot{x}+4y=0,\\
      &x(0)=5,\\
      &y(0)=\dot{y}(0)=\dot{x}(0)=0.
   \end{aligned}
   \right.
   \]
    \begin{enumerate}
       \item (1+) [2 points] Via a reduction of order, convert this system to a system of first-order equations.
       \item (1+) [2 points] Solve the system discovered in part (a).
    \end{enumerate}

   \item (2) [3 points] Verify that the characteristic polynomial of the matrix \(A\) of the system 
   \[
   \vec{x}'=\begin{pmatrix}
      3&-4&1&0\\4&3&0&1\\0&0&3&-4\\0&0&4&3
   \end{pmatrix}\vec{x}
   \]
   is \(f_A(\lambda)=(\lambda^2-6\lambda+25)^2\) therefore showing that \(A\) has a repeated complex conjugate pair \(3\pm 4i\) of eigenvalues.\par
   Verify that the complex vectors 
   \[
   \vec{x}_{\lambda,1}=(1,i,0,0),\quad\vec{x}_{\lambda,2}=(0,0,1,i),
   \]
   are associated with the eigenvalue \(\lambda=3-4i\) and then calculate the real and imaginary parts of the complex-valued solutions 
   \[
   \vec{x}_{\lambda,1}e^{\lambda t},\quad\text{and}\quad(\vec{x}_{\lambda,1}t+\vec{x}_{\lambda,2})e^{\lambda t}
   \]
   to find the four real valued solutions of \(\vec{x}'=A\vec{x}\).

    \item (2+) [4 points] Solve the system of differential equations given by 
   \[
   \left\lbrace
   \begin{aligned}
      &\dot{x}(t)+\dot{y}(t)=z(t)\\
      &\dot{y}(t)+\dot{z}(t)=x(t)\\
      &\dot{z}(t)+\dot{x}(t)=y(t)\\
   \end{aligned}
   \right.
   \]
   with initial condition \((x(0),y(0),z(0))=(1,1,1)\).

    \item (1-) [1 point] Transform the following systems of ODEs into linear differential equations.
    \[
   \left\lbrace
   \begin{aligned}
      &\dot{x}(t)=y(t),\\
      &\dot{y}(t)=t+t^2-e^{-4t},\\
   \end{aligned}
   \right.
   \quad\text{and}\quad
   \left\lbrace
   \begin{aligned}
      &\dot{x}(t)=y(t),\\
      &\dot{y}(t)=z(t),\\
      &\dot{z}(t)=\int_0^t\frac{\sin(u)}{u}\dd u.\\
   \end{aligned}
   \right.
   \]

   \item (2) [2 points] Reduce the order of the following higher order ODE by turning them into systems of differential equations.
   \[
   x'''+xx''=\cos(t),\quad\text{and}\quad (y''')^2-\sin(x)y'=y\cos(x).
   \]


\end{enumerate}

\end{document}